\section{Zielsetzung und Forschungsfragen}

Ziel dieser Arbeit ist die Untersuchung, unter welchen Voraussetzungen ERP-Systeme in kleinen und mittleren Unternehmen als technologische 
Basis digitaler Transformation wirksam werden können. Im Mittelpunkt steht dabei nicht ausschließlich die Einführung von ERP-Systemen, 
sondern insbesondere deren organisatorische und strategische Wirksamkeit im Kontext digitaler Transformationsprozesse.

Die Arbeit verfolgt das Ziel, zentrale Voraussetzungen für die Wirksamkeit von ERP-Systemen in KMU systematisch 
zu identifizieren und in einem konzeptionellen ERP-KMU-Fit-Modell zusammenzuführen. Das Modell soll dazu beitragen, unterschiedliche 
organisatorische, technologische und strategische Einflussfaktoren strukturiert zu betrachten und deren Bedeutung für die 
Transformationsfähigkeit von ERP-Systemen in KMU einzuordnen.

\noindent Daraus ergibt sich folgende Forschungsfrage:

\begin{quote}
\glqq Welche Voraussetzungen müssen KMU erfüllen, damit ERP-Systeme als technologische Basis digitaler Transformation wirksam werden können?\grqq{}.
\end{quote}

\noindent Zur Beantwortung dieser Forschungsfrage werden literaturbasierte Erfolgsbedingungen analysiert, theoretisch eingeordnet und anschließend 
innerhalb eines konzeptionellen Modells zusammengeführt.